% Generated draft: requires mathematical and source review.
\documentclass[11pt]{article}
\usepackage[margin=25mm]{geometry}
\usepackage{fontspec}
\setmainfont{Latin Modern Roman}
\usepackage{amsmath,amssymb,mathtools,unicode-math}
\setmathfont{Latin Modern Math}
\usepackage{enumitem,needspace,xurl,hyperref}
\setlength{\parindent}{0pt}
\setlength{\parskip}{5pt}
\setlength{\emergencystretch}{3em}
\newcommand{\N}{\mathbb{N}}
\newcommand{\Z}{\mathbb{Z}}
\newcommand{\Q}{\mathbb{Q}}
\newcommand{\R}{\mathbb{R}}
\newcommand{\C}{\mathbb{C}}
\newcommand{\F}{\mathbb{F}}
\newcommand{\A}{\mathbb{A}}
\newcommand{\PP}{\mathbb{P}}
\newcommand{\E}{\mathbb{E}}
\newcommand{\eps}{\varepsilon}
\newcommand{\dd}{\,\mathrm{d}}
\begin{document}
\section*{UnsolvedMath ID: 30005496}
\textbf{Source problem code:} catalogue entry
\subsection{Definitions and mathematical statement}
The supplied question is
\[
\text{For every rational log-canonical singularity }(X;x),\quad
\text{is }\pi_1(X;x)\text{ virtually nilpotent?}
\]
A group \(G\) is \emph{virtually nilpotent} if there is a subgroup \(H\leq G\) with finite index \([G:H]<\infty\) such that its lower central series, defined by
\[
\gamma_1(H)=H,\qquad \gamma_{j+1}(H)=[H,\gamma_j(H)],
\]
satisfies \(\gamma_{c+1}(H)=\{1\}\) for some integer \(c\geq 1\). Thus the requested conclusion is
\[
\forall (X;x)\text{ rational and log-canonical},\quad
\exists H\leq\pi_1(X;x)\ \exists c\geq1:\quad
[\pi_1(X;x):H]<\infty,\qquad \gamma_{c+1}(H)=\{1\}.
\]
No uniform bound on the index or the nilpotency class is asserted.

Under the conventional interpretation as a normal complex variety germ, rationality means that, for a resolution \(f:Y\to V\) of a representative neighborhood of \(x\), the higher direct images \(R^i f_*\mathcal O_Y\) vanish near \(x\) for every \(i>0\). Under the boundary-free interpretation, log canonicity means that \(K_V\) is \(\Q\)-Cartier and every prime divisor \(E\) over \(V\) with center near \(x\) has discrepancy \(a(E,V,0)\geq-1\); for exceptional divisors on a resolution these discrepancies occur in
\[
K_Y=f^*K_V+\sum_{E\text{ exceptional}}a(E,V,0)E.
\]
These definitions describe a conventional interpretation, since the supplied record does not explicitly specify the ground field, the algebraic or analytic category, or whether a boundary is involved.

The notation \(\pi_1(X;x)\) is also unresolved in the supplied record. In a complex setting, the fundamental group of a sufficiently small punctured neighborhood, \(\pi_1(U\setminus\{x\},u)\), and that of its regular locus, \(\pi_1(U\cap X_{\mathrm{reg}},u)\), are distinct possible targets when the singularity is not isolated. The record does not specify either convention, does not assert isolatedness, and does not explicitly distinguish a topological local group from an algebraic local fundamental group. The displayed question therefore retains the source notation pending verification of its intended definition.
\subsection{Short English statement}
The supplied question asks whether every rational log-canonical singularity has a local fundamental group containing a nilpotent subgroup of finite index. The meaning of the local fundamental group and the ambient geometric conventions require confirmation from the report's surrounding discussion. The supplied open status is dated curation, not an independently verified conclusion.
\subsection{Sources}
\begin{itemize}
\item Supplied local UnsolvedMath registry record, permanent ID \(30005496\), title \emph{Virtual Nilpotence of Rational Log-Canonical Local Groups}, exact source problem code \texttt{catalogue entry}, in the collection \emph{Oberwolfach Reports Open Problems}. The record gives proposed year 2023 but names no individual proposer.
\item Supplied report citation: \emph{Komplexe Analysis -- Differential and Algebraic methods in Kähler spaces} (2023), Oberwolfach Reports, pp.~965--1030, DOI 10.4171/owr/2023/18. \url{https://doi.org/10.4171/owr/2023/18}. The supplied original excerpt asks whether \(\pi_1(X;x)\) is virtually nilpotent and includes a page marker for p.~1009; the rational and log-canonical hypotheses appear in the supplied clean statement rather than in that excerpt.
\item Supplied literature source: \emph{Low-dimensional lc singularities}. \url{https://arxiv.org/abs/2302.11790}.
\item Supplied literature source: \emph{OWR 2023 report}. \url{https://ems.press/content/serial-article-files/47013}.
\item Literature assessment dated exactly \texttt{2026-08-22}: the supplied curation labels the question open, reports that no general proof was found, and says that the report presents it as an open question. This is supplied dated curation derived from the OWR record, not independently verified literature status.
\item The registry labels the statement \texttt{corrected\_verified} and reports that it was checked against the official Oberwolfach report with its hypotheses, notation, and target preserved. This verification claim is supplied metadata; no independent source consultation is claimed here.
\end{itemize}
\end{document}
